% ======================================================================
%  PLANTILLA — Informe de una entrega "Programa N"
%  Proyecto Integrador: Calculadora de Álgebra Lineal — UAM (MTM0120)
%
%  Este archivo SÍ se versiona (portada + estructura reutilizables).
%  Los datos personales (nombres legales de los integrantes, grupo, docente,
%  fecha) van en un `datos.tex` que NO se versiona. Ver `informes/README.md`.
%
%  Uso:
%    1. cp -r informes/plantilla informes/Informe_Programa_N_Grupo_7
%    2. cd informes/Informe_Programa_N_Grupo_7
%    3. cp datos.example.tex datos.tex   &&   editar datos.tex
%    4. poner img/uam-logo.png y las capturas img/casoN-*.png
%    5. pdflatex informe.tex && pdflatex informe.tex   (dos veces, por el índice)
%    6. renombrar informe.pdf -> "Informe_Programa N_Grupo 7.pdf"
% ======================================================================
\documentclass[11pt,a4paper]{article}

\usepackage[T1]{fontenc}
\usepackage[utf8]{inputenc}
\usepackage{lmodern}
\usepackage[spanish,es-noindentfirst]{babel}
\usepackage[a4paper,margin=2.5cm]{geometry}
\usepackage{amsmath,amssymb}
\usepackage{graphicx}
\usepackage{booktabs}
\usepackage{enumitem}
\usepackage{xcolor}
\usepackage{listings}
\usepackage[hidelinks]{hyperref}
\usepackage{parskip}

\graphicspath{{img/}}

% --- Datos de la entrega -------------------------------------------------
% Valores por defecto (para que la plantilla compile sin datos.tex).
% Los reales viven en datos.tex, que NO se versiona.
\newcommand{\programaTitulo}{Programa N: [t\'itulo de la tarea]}
\newcommand{\grupoNum}{7}
\newcommand{\docente}{Jos\'e Andr\'es Mungu\'ia Cort\'ez}
\newcommand{\ciudadFecha}{Managua, \today}
\newcommand{\integrantes}{[integrantes: ver datos.tex]}

\IfFileExists{datos.tex}{\input{datos.tex}}{%
  \typeout{*** AVISO: falta datos.tex; usando valores por defecto (ver informes/README.md) ***}%
}
% ----------------------------------------------------------------------

% --- Caja de pseudocódigo (imita los recuadros oscuros del informe 1) -
\definecolor{codebg}{HTML}{2B2B2B}
\definecolor{codefg}{HTML}{E6E6E6}
\lstdefinestyle{pseudo}{
  backgroundcolor=\color{codebg}, basicstyle=\ttfamily\small\color{codefg},
  breaklines=true, columns=fullflexible, keepspaces=true,
  frame=single, framerule=0pt, framesep=8pt, xleftmargin=8pt, xrightmargin=8pt,
  aboveskip=1em, belowskip=1em, literate={->}{{$\rightarrow$}}1 {≤}{{$\leq$}}1 {≠}{{$\neq$}}1
}
\lstset{style=pseudo}

% --- Una captura (o recuadro [PENDIENTE] si aún no se tomó) ----------
\newcommand{\captura}[2]{%
  \begin{center}
  \IfFileExists{img/#1}{%
    \fbox{\includegraphics[width=0.96\linewidth]{#1}}\\[0.3em]
    {\footnotesize #2}%
  }{%
    \fbox{\parbox{0.9\linewidth}{\centering\vspace{1.2em}
      \textbf{[PENDIENTE: captura de pantalla]}\\[0.4em]
      \small #2\\[0.2em]
      Guardar como \texttt{img/#1}
      (\texttt{informes/plantilla/generar-capturas.py}).
      \vspace{1.2em}}}%
  }
  \end{center}%
}
% --- Las dos capturas de un caso: pantalla de proceso + de vector ----
\newcommand{\capturaCaso}[2]{%
  \captura{#1-proceso.png}{Pantalla \emph{Proceso y resultado}: #2}%
  \captura{#1-vectorial.png}{Pantalla \emph{Soluci\'on en notaci\'on vectorial}.}%
}

\begin{document}

% ===================== PORTADA =======================================
\begin{titlepage}
\centering
{\large\bfseries UNIVERSIDAD AMERICANA\par}
\vspace{0.4em}
\rule{\linewidth}{0.4pt}\par
\vspace{1.5em}

\IfFileExists{img/uam-logo.png}{\includegraphics[width=0.42\linewidth]{uam-logo.png}}{%
  \fbox{\parbox{0.42\linewidth}{\centering\vspace{3em}\small [logo UAM: \texttt{img/uam-logo.png}]\vspace{3em}}}}
\par\vspace{1.5em}

{Asignatura: \'Algebra Lineal (MTM0120)\par}
{Facultad de Ingenier\'ia y Arquitectura (FIA)\par}
\vspace{1em}
\rule{\linewidth}{0.4pt}\par
\vspace{1.2em}

{\itshape\bfseries Proyecto Integrador: Calculadora de \'Algebra Lineal\par}
\vspace{0.3em}
{\itshape\bfseries \programaTitulo\par}
\vspace{1.2em}
\rule{\linewidth}{0.4pt}\par
\vspace{2em}

{\bfseries Integrantes:\par}
\vspace{0.3em}
{\integrantes\par}
\vspace{3em}

{\bfseries Grupo:} \grupoNum\par
\vspace{2em}
{\bfseries Docente:} \docente\par
\vspace{0.5em}
{Carrera: Ingenier\'ia en Sistemas y Tecnolog\'ias de la Informaci\'on\par}

\vfill
{\ciudadFecha\par}
\end{titlepage}

% ===================== ÍNDICE ========================================
\tableofcontents
\newpage

% ===================== EXPLICACIÓN DEL ALGORITMO =====================
\section{Explicaci\'on del algoritmo}

\subsection{Contexto y ubicaci\'on en el repositorio}

El programa transforma una matriz aumentada $[\,A \mid b\,]$ a su \textbf{Forma
Escalonada Reducida por Filas} (RREF) mediante el m\'etodo de \textbf{Gauss--Jordan},
e informa de forma expl\'icita las \textbf{columnas pivote}, las \textbf{variables
b\'asicas} y las \textbf{variables libres}, y la soluci\'on (\'unica o su estructura
general). Todo se implementa con Python est\'andar: listas anidadas, condicionales,
bucles y funciones propias; sin NumPy, SciPy ni funciones de \'algebra lineal de
\texttt{math}. Los m\'odulos de c\'alculo no hacen \texttt{input()}, \texttt{print()}
ni acceso a archivos --- son funciones puras verificables por separado.

En el repositorio del curso (\texttt{alg-lineal-I}) la l\'ogica vive en el n\'ucleo
de c\'alculo:
\begin{itemize}[nosep]
  \item \texttt{gauss.py} --- \texttt{escalonar()}: reducci\'on a forma escalonada
        (REF) con pivoteo parcial; devuelve la lista \texttt{columnas\_pivote}.
        Adem\'as \texttt{rango\_matriz()}, \texttt{verificar\_rango\_nulidad()},
        \texttt{clasificar()}.
  \item \texttt{gauss\_jordan.py} --- \texttt{reducir\_a\_escalonada\_reducida()}:
        contin\'ua de REF a RREF (normaliza cada pivote a 1 y hace ceros
        \emph{encima} del pivote); \texttt{solucion\_parametrica()} y
        \texttt{solucion\_general\_vectorial()}.
\end{itemize}

\subsection{Fundamentos matem\'aticos}

La reducci\'on usa \'unicamente \textbf{operaciones elementales de fila}, que no
alteran el conjunto soluci\'on:
\begin{itemize}[nosep]
  \item Intercambiar dos filas ($F_i \leftrightarrow F_j$).
  \item Multiplicar una fila por un escalar no nulo ($F_i \leftarrow \tfrac{1}{p}\,F_i$).
  \item Restarle a una fila un m\'ultiplo de otra ($F_i \leftarrow F_i - \lambda\,F_j$).
\end{itemize}

\paragraph{Pivoteo parcial.} En cada columna se elige como pivote, entre la fila
actual y las inferiores, la de mayor valor absoluto. Esto evita dividir por cero
cuando el elemento en la diagonal es $0$ y reduce el error de redondeo del punto
flotante. Las comparaciones con cero usan una tolerancia $\varepsilon = 10^{-9}$
(\texttt{valor\_casi\_cero}), nunca \texttt{== 0}.

\subsection{Fase 1: reducci\'on a forma escalonada (\texttt{escalonar})}

Se recorren las columnas de izquierda a derecha manteniendo un puntero
\texttt{fila\_actual}. Para la columna $j$:

\begin{enumerate}[nosep]
  \item \textbf{Selecci\'on de pivote:} buscar, de \texttt{fila\_actual} hacia
        abajo, la fila con $|a_{rj}|$ m\'aximo.
  \item Si ese m\'aximo es $\approx 0$, la columna $j$ no tiene pivote: ser\'a una
        \textbf{variable libre}. Se pasa a la siguiente columna sin avanzar
        \texttt{fila\_actual}.
  \item \textbf{Intercambio} de la fila del pivote con \texttt{fila\_actual} (si
        difieren); se registra el paso $F_i \leftrightarrow F_k$.
  \item \textbf{Ceros debajo del pivote:} para cada fila inferior $r$,
        $\lambda = a_{rj}/p$ y $F_r \leftarrow F_r - \lambda\,F_{\text{actual}}$.
  \item Registrar $j$ en \texttt{columnas\_pivote} y avanzar \texttt{fila\_actual}.
\end{enumerate}

\begin{lstlisting}
Para cada columna j de 0 a n-1:
    fila_max = argmax(|a[r][j]|)  para r >= fila_actual
    si |a[fila_max][j]| ~ 0:  continuar        # columna sin pivote -> variable libre
    intercambiar(fila_actual, fila_max)
    p = a[fila_actual][j]
    Para cada fila r > fila_actual:
        lambda = a[r][j] / p
        a[r] = a[r] - lambda * a[fila_actual]  # anula a[r][j]
    columnas_pivote.append(j);  fila_actual += 1
devolver columnas_pivote
\end{lstlisting}

\noindent\textbf{Resultado:} matriz en REF y \texttt{columnas\_pivote} = \'indices
(base 0) de las columnas con pivote. Su longitud es el \textbf{rango} de $A$.

\subsection{Fase 2: ceros por encima del pivote (\texttt{reducir\_a\_escalonada\_reducida})}

Partiendo de la REF, se recorren los pivotes \textbf{de abajo hacia arriba}. Para el
pivote $i$ en la columna \texttt{col}:

\begin{enumerate}[nosep]
  \item \textbf{Normalizar a 1:} si $a_{i,\text{col}} \neq 1$, dividir toda la fila
        $i$ (desde \texttt{col}) entre el pivote. Paso: $F_i \leftarrow F_i / p$.
  \item \textbf{Ceros arriba:} para cada fila $r < i$, tomar
        $\lambda = a_{r,\text{col}}$; si $\lambda \not\approx 0$,
        $F_r \leftarrow F_r - \lambda\,F_i$ y fijar $a_{r,\text{col}} = 0$
        expl\'icitamente (evita residuos de punto flotante).
\end{enumerate}

\begin{lstlisting}
Para i de (num_pivotes - 1) hasta 0:          # de abajo hacia arriba
    col = columnas_pivote[i];  p = a[i][col]
    si p != 1:  a[i] = a[i] / p               # F_i <- F_i / p
    Para cada fila r < i:
        lambda = a[r][col]
        si lambda ~ 0:  continuar
        a[r] = a[r] - lambda * a[i]           # F_r <- F_r - lambda * F_i
        a[r][col] = 0.0
\end{lstlisting}

\noindent\textbf{Resultado (RREF):} cada pivote vale $1$ y es el \'unico elemento no
nulo de su columna; toda columna pivote es un vector can\'onico $e_k$.

\subsection{Identificaci\'on de columnas pivote y variables libres}

\begin{itemize}[nosep]
  \item \textbf{Columnas pivote} $=$ \texttt{columnas\_pivote} (Fase 1). En el
        reporte se muestran en base 1: \emph{``Las columnas pivote son: 1, 2 y 4''}.
  \item \textbf{Variables b\'asicas} $=$ las asociadas a una columna pivote.
  \item \textbf{Variables libres} $=$
        \texttt{[v for v in range(n) if v not in columnas\_pivote]}.
  \item \textbf{Rango} $= |\texttt{columnas\_pivote}|$;\quad
        \textbf{Nulidad} $= n - \text{rango}$;\quad
        siempre $\operatorname{Rango}(A) + \operatorname{Nulidad}(A) = n$.
\end{itemize}

\subsection{Clasificaci\'on del sistema}

Sobre la matriz reducida (criterio de Rouch\'e--Frobenius):
\begin{itemize}[nosep]
  \item \textbf{Inconsistente:} existe una fila $[\,0\ \cdots\ 0 \mid c\,]$ con
        $c \neq 0$.
  \item \textbf{Consistente determinado:} no hay fila de ese tipo y
        $\text{rango} = n$ (todas las variables son b\'asicas) $\Rightarrow$
        soluci\'on \'unica.
  \item \textbf{Consistente indeterminado:} no hay fila de ese tipo y
        $\text{rango} < n$ $\Rightarrow$ infinitas soluciones (hay variables
        libres).
\end{itemize}

\subsection{Soluci\'on a partir de la RREF}

\paragraph{Determinado.} Con la RREF, la columna $b$ ya es la soluci\'on: la fila
del pivote de $x_k$ dice $x_k = b'_k$.

\paragraph{Indeterminado.} \texttt{solucion\_parametrica()} lee cada fila pivote de
la RREF y despeja su variable b\'asica en funci\'on de las libres:
\[
  x_{\text{piv}} \;=\; b'_i \;-\!\!\sum_{v\ \text{libre}} a'_{i,v}\,x_v .
\]
\texttt{solucion\_general\_vectorial()} arma entonces
\[
  \mathbf{x} \;=\; \mathbf{x}_p \;+\; t_1 \mathbf{v}_1 \;+\; \cdots \;+\; t_k \mathbf{v}_k ,
\]
donde $\mathbf{x}_p$ se obtiene fijando todas las libres a $0$, y cada
$\mathbf{v}_j$ fijando $x_{\text{libre}_j}=1$ y el resto de libres a $0$ y restando
$\mathbf{x}_p$ ($k = $ nulidad).

% ===================== CASOS DE PRUEBA ==============================
\section{Casos de prueba (capturas de pantalla)}

Ejecuci\'on en la interfaz gr\'afica (\texttt{uv run aqua-gauss}). Cada caso muestra
la matriz llevada a RREF, las columnas pivote, la clasificaci\'on de variables y la
soluci\'on.

\subsection{Sistema con soluci\'on \'unica}

\[
\begin{array}{rcrcrcr}
 x_1 &+& x_2 &+& x_3 &=& 6\\
     & & 2x_2 &+& 5x_3 &=& -4\\
 2x_1 &+& 5x_2 &-& x_3 &=& 27
\end{array}
\qquad
[\,A\mid b\,]=\left[\begin{array}{ccc|c}
1&1&1&6\\ 0&2&5&-4\\ 2&5&-1&27
\end{array}\right]
\]

RREF esperada $\left[\begin{smallmatrix}1&0&0&5\\0&1&0&3\\0&0&1&-2\end{smallmatrix}\right]$;
columnas pivote: 1, 2 y 3; todas las variables son b\'asicas (no hay libres);
$\text{rango}=3=n$. Soluci\'on \'unica $x_1=5,\ x_2=3,\ x_3=-2$.

\capturaCaso{caso1}{el \'ultimo paso es la forma escalonada final; abajo, la
soluci\'on \'unica $(5,3,-2)$ y su comprobaci\'on. Columnas pivote 1, 2 y 3
(todas b\'asicas), $\text{rango}=3$.}

\subsection{Sistema con variables libres (infinitas soluciones)}

\[
\begin{array}{rcrcrcr}
 x_1 &+& x_2 &+& x_3 &=& 6\\
 2x_1 &+& 2x_2 &+& 2x_3 &=& 12\\
 x_1 &-& x_2 & & &=& 0
\end{array}
\qquad
[\,A\mid b\,]=\left[\begin{array}{ccc|c}
1&1&1&6\\ 2&2&2&12\\ 1&-1&0&0
\end{array}\right]
\]

La fila $2$ es $2\times$ la fila $1$: se anula al escalonar. RREF esperada
$\left[\begin{smallmatrix}1&0&\tfrac12&3\\[2pt]0&1&\tfrac12&3\\[2pt]0&0&0&0\end{smallmatrix}\right]$;
columnas pivote: 1 y 2; \textbf{variable libre:} $x_3$; $\text{rango}=2$,
$\text{nulidad}=1$. Soluci\'on general
\[
\mathbf{x}=\begin{pmatrix}3\\3\\0\end{pmatrix}
      + t_1\begin{pmatrix}-\tfrac12\\[2pt]-\tfrac12\\[2pt]1\end{pmatrix}.
\]

\capturaCaso{caso2}{el \'ultimo paso (10 de 10) es la \textbf{RREF}
$\left[\begin{smallmatrix}1&0&\tfrac12&3\\0&1&\tfrac12&3\\0&0&0&0\end{smallmatrix}\right]$;
columnas pivote 1 y 2, $x_3$ libre, $\text{rango}=2$, $\text{nulidad}=1$.}

\subsection{Sistema sin soluci\'on (inconsistente)}

\[
\begin{array}{rcrcr}
 x_1 &+& x_2 &=& 2\\
 x_1 &+& x_2 &=& 5
\end{array}
\qquad
[\,A\mid b\,]=\left[\begin{array}{cc|c}
1&1&2\\ 1&1&5
\end{array}\right]
\]

Al restar $F_2 \leftarrow F_2 - F_1$ queda $[\,0\ \ 0 \mid 3\,]$: una ecuaci\'on
$0 = 3$. El sistema es \textbf{inconsistente}; no hay soluci\'on ni forma
vectorial.

\capturaCaso{caso3}{la forma escalonada final tiene la fila $[\,0\ \ 0\mid 3\,]$
(``$0 = 3$''); estado en rojo ``Sistema inconsistente''. No hay soluci\'on ni
forma vectorial.}

% ===================== CONCLUSIÓN ===================================
\section{Conclusi\'on}

El programa cumple los requerimientos de la Tarea 2: entrada de $[\,A\mid b\,]$,
reducci\'on completa a RREF por Gauss--Jordan (pivotes a $1$, ceros encima y
debajo), detecci\'on autom\'atica de columnas pivote, distinci\'on entre variables
b\'asicas y libres, y despliegue de la soluci\'on \'unica o de la estructura de la
soluci\'on general --- todo con Python est\'andar, sin librer\'ias de \'algebra
lineal. La fase de ceros por encima del pivote reutiliza \texttt{columnas\_pivote}
de la fase de escalonado, de modo que identificar los pivotes y construir la RREF
son un mismo recorrido de la matriz.

\end{document}
